% Options for packages loaded elsewhere
\PassOptionsToPackage{unicode}{hyperref}
\PassOptionsToPackage{hyphens}{url}
\PassOptionsToPackage{dvipsnames,svgnames,x11names}{xcolor}
\documentclass[
  spanish,
  10pt,
  a4paper,
]{article}
\usepackage{xcolor}
\usepackage[top=2.2cm,bottom=2.2cm,left=2.1cm,right=2.1cm]{geometry}
\usepackage{amsmath,amssymb}
\setcounter{secnumdepth}{5}
\usepackage{iftex}
\ifPDFTeX
  \usepackage[T1]{fontenc}
  \usepackage[utf8]{inputenc}
  \usepackage{textcomp} % provide euro and other symbols
\else % if luatex or xetex
  \usepackage{unicode-math} % this also loads fontspec
  \defaultfontfeatures{Scale=MatchLowercase}
  \defaultfontfeatures[\rmfamily]{Ligatures=TeX,Scale=1}
\fi
\usepackage{lmodern}
\ifPDFTeX\else
  % xetex/luatex font selection
\fi
% Use upquote if available, for straight quotes in verbatim environments
\IfFileExists{upquote.sty}{\usepackage{upquote}}{}
\IfFileExists{microtype.sty}{% use microtype if available
  \usepackage[]{microtype}
  \UseMicrotypeSet[protrusion]{basicmath} % disable protrusion for tt fonts
}{}
\makeatletter
\@ifundefined{KOMAClassName}{% if non-KOMA class
  \IfFileExists{parskip.sty}{%
    \usepackage{parskip}
  }{% else
    \setlength{\parindent}{0pt}
    \setlength{\parskip}{6pt plus 2pt minus 1pt}}
}{% if KOMA class
  \KOMAoptions{parskip=half}}
\makeatother
\usepackage{color}
\usepackage{fancyvrb}
\newcommand{\VerbBar}{|}
\newcommand{\VERB}{\Verb[commandchars=\\\{\}]}
\DefineVerbatimEnvironment{Highlighting}{Verbatim}{commandchars=\\\{\}}
% Add ',fontsize=\small' for more characters per line
\newenvironment{Shaded}{}{}
\newcommand{\AlertTok}[1]{\textcolor[rgb]{1.00,0.00,0.00}{\textbf{#1}}}
\newcommand{\AnnotationTok}[1]{\textcolor[rgb]{0.38,0.63,0.69}{\textbf{\textit{#1}}}}
\newcommand{\AttributeTok}[1]{\textcolor[rgb]{0.49,0.56,0.16}{#1}}
\newcommand{\BaseNTok}[1]{\textcolor[rgb]{0.25,0.63,0.44}{#1}}
\newcommand{\BuiltInTok}[1]{\textcolor[rgb]{0.00,0.50,0.00}{#1}}
\newcommand{\CharTok}[1]{\textcolor[rgb]{0.25,0.44,0.63}{#1}}
\newcommand{\CommentTok}[1]{\textcolor[rgb]{0.38,0.63,0.69}{\textit{#1}}}
\newcommand{\CommentVarTok}[1]{\textcolor[rgb]{0.38,0.63,0.69}{\textbf{\textit{#1}}}}
\newcommand{\ConstantTok}[1]{\textcolor[rgb]{0.53,0.00,0.00}{#1}}
\newcommand{\ControlFlowTok}[1]{\textcolor[rgb]{0.00,0.44,0.13}{\textbf{#1}}}
\newcommand{\DataTypeTok}[1]{\textcolor[rgb]{0.56,0.13,0.00}{#1}}
\newcommand{\DecValTok}[1]{\textcolor[rgb]{0.25,0.63,0.44}{#1}}
\newcommand{\DocumentationTok}[1]{\textcolor[rgb]{0.73,0.13,0.13}{\textit{#1}}}
\newcommand{\ErrorTok}[1]{\textcolor[rgb]{1.00,0.00,0.00}{\textbf{#1}}}
\newcommand{\ExtensionTok}[1]{#1}
\newcommand{\FloatTok}[1]{\textcolor[rgb]{0.25,0.63,0.44}{#1}}
\newcommand{\FunctionTok}[1]{\textcolor[rgb]{0.02,0.16,0.49}{#1}}
\newcommand{\ImportTok}[1]{\textcolor[rgb]{0.00,0.50,0.00}{\textbf{#1}}}
\newcommand{\InformationTok}[1]{\textcolor[rgb]{0.38,0.63,0.69}{\textbf{\textit{#1}}}}
\newcommand{\KeywordTok}[1]{\textcolor[rgb]{0.00,0.44,0.13}{\textbf{#1}}}
\newcommand{\NormalTok}[1]{#1}
\newcommand{\OperatorTok}[1]{\textcolor[rgb]{0.40,0.40,0.40}{#1}}
\newcommand{\OtherTok}[1]{\textcolor[rgb]{0.00,0.44,0.13}{#1}}
\newcommand{\PreprocessorTok}[1]{\textcolor[rgb]{0.74,0.48,0.00}{#1}}
\newcommand{\RegionMarkerTok}[1]{#1}
\newcommand{\SpecialCharTok}[1]{\textcolor[rgb]{0.25,0.44,0.63}{#1}}
\newcommand{\SpecialStringTok}[1]{\textcolor[rgb]{0.73,0.40,0.53}{#1}}
\newcommand{\StringTok}[1]{\textcolor[rgb]{0.25,0.44,0.63}{#1}}
\newcommand{\VariableTok}[1]{\textcolor[rgb]{0.10,0.09,0.49}{#1}}
\newcommand{\VerbatimStringTok}[1]{\textcolor[rgb]{0.25,0.44,0.63}{#1}}
\newcommand{\WarningTok}[1]{\textcolor[rgb]{0.38,0.63,0.69}{\textbf{\textit{#1}}}}
\usepackage{longtable,booktabs,array}
\newcounter{none} % for unnumbered tables
\usepackage{calc} % for calculating minipage widths
% Correct order of tables after \paragraph or \subparagraph
\usepackage{etoolbox}
\makeatletter
\patchcmd\longtable{\par}{\if@noskipsec\mbox{}\fi\par}{}{}
\makeatother
% Allow footnotes in longtable head/foot
\IfFileExists{footnotehyper.sty}{\usepackage{footnotehyper}}{\usepackage{footnote}}
\makesavenoteenv{longtable}
\ifLuaTeX
\usepackage[bidi=basic,shorthands=off]{babel}
\else
\usepackage[bidi=default,shorthands=off]{babel}
\fi
\ifLuaTeX
  \usepackage{selnolig} % disable illegal ligatures
\fi
\setlength{\emergencystretch}{3em} % prevent overfull lines
\providecommand{\tightlist}{%
  \setlength{\itemsep}{0pt}\setlength{\parskip}{0pt}}
% Shared visual layer for the generated Kaleido FlowTwin reports.
\usepackage{helvet}
\renewcommand{\familydefault}{\sfdefault}
\usepackage{microtype}
\usepackage{xcolor}
\usepackage{titlesec}
\usepackage{fancyhdr}
\usepackage{enumitem}
\usepackage{booktabs}
\usepackage{xurl}
\usepackage{lastpage}
\usepackage{etoolbox}

\definecolor{KaleidoNavy}{HTML}{0A2342}
\definecolor{KaleidoBlue}{HTML}{176B87}
\definecolor{KaleidoCyan}{HTML}{3BB7C4}
\definecolor{KaleidoMint}{HTML}{DDF5F1}
\definecolor{KaleidoFog}{HTML}{F2F5F7}
\definecolor{KaleidoSlate}{HTML}{5D6A7A}
\definecolor{KaleidoCoral}{HTML}{F26B5B}

\titleformat{\section}
  {\Large\bfseries\color{KaleidoNavy}}
  {\thesection}{0.75em}{}[\vspace{0.2em}\color{KaleidoCyan}\titlerule]
\titleformat{\subsection}
  {\large\bfseries\color{KaleidoBlue}}
  {\thesubsection}{0.65em}{}
\titleformat{\subsubsection}
  {\normalsize\bfseries\color{KaleidoNavy}}
  {\thesubsubsection}{0.55em}{}
\titlespacing*{\section}{0pt}{2.2em}{1em}
\titlespacing*{\subsection}{0pt}{1.6em}{0.65em}

\pagestyle{fancy}
\fancyhf{}
\fancyhead[L]{\small\bfseries\color{KaleidoNavy}KALEIDO FLOWTWIN}
\fancyhead[R]{\small\color{KaleidoSlate}DOSSIER TECNICO}
\fancyfoot[L]{\scriptsize\color{KaleidoSlate}EVOCON Solutions GmbH}
\fancyfoot[R]{\scriptsize\color{KaleidoSlate}\thepage/\pageref*{LastPage}}
\renewcommand{\headrulewidth}{0.6pt}
\renewcommand{\headrule}{\hbox to\headwidth{\color{KaleidoCyan}\leaders\hrule height \headrulewidth\hfill}}
\renewcommand{\footrulewidth}{0pt}
\setlength{\headheight}{20pt}

\setlist[itemize]{leftmargin=1.4em,itemsep=0.25em,topsep=0.35em}
\setlist[enumerate]{leftmargin=1.7em,itemsep=0.25em,topsep=0.35em}
\setlength{\parindent}{0pt}
\setlength{\parskip}{0.55em}
\setlength{\emergencystretch}{3em}
\renewcommand{\arraystretch}{1.18}

\AtBeginEnvironment{quote}{%
  \color{KaleidoNavy}\itshape
  \setlength{\leftskip}{1em}
  \setlength{\rightskip}{1em}
}

\makeatletter
\renewcommand{\maketitle}{%
  \hypersetup{pageanchor=false}
  \begin{titlepage}
    \pagecolor{KaleidoNavy}
    \color{white}
    \vspace*{1.4cm}
    {\color{KaleidoCyan}\rule{2.4cm}{2.5pt}\par}
    \vspace{1.1cm}
    {\Huge\bfseries\@title\par}
    \vspace{0.65cm}
    {\Large\color{KaleidoMint}Kaleido FlowTwin\par}
    \vfill
    {\large Álvaro Schwiedop Souto\par}
    \vspace{0.15cm}
    {\normalsize\color{KaleidoMint}Industrial AI \& Predictive Quality Lead\par}
    \vspace{0.15cm}
    {\normalsize\color{KaleidoMint}EVOCON Solutions GmbH\par}
    \vspace{0.25cm}
    {\normalsize\color{KaleidoMint}\@date\par}
  \end{titlepage}
  \hypersetup{pageanchor=true}
  \nopagecolor
  \color{black}
}
\makeatother
\usepackage{bookmark}
\IfFileExists{xurl.sty}{\usepackage{xurl}}{} % add URL line breaks if available
\urlstyle{same}
% fallback for those not using the hyperref driver hyperxmp:
\makeatletter
\@ifundefined{xmpquote}{\newcommand{\xmpquote}[1]{#1}}{}
\makeatother
\hypersetup{
  pdftitle={Guion de reunion: Kaleido FlowTwin},
  pdflang={es},
  colorlinks=true,
  linkcolor={KaleidoBlue},
  filecolor={Maroon},
  citecolor={Blue},
  urlcolor={KaleidoBlue},
  pdfcreator={LaTeX via pandoc}}

\title{Guion de reunion: Kaleido FlowTwin}
\author{}
\date{15 julio 2026}

\begin{document}
\maketitle

{
\setcounter{tocdepth}{3}
\tableofcontents
}
\section{Guion de reunion: Kaleido
FlowTwin}\label{guion-de-reunion-kaleido-flowtwin}

Documento de preparacion para una conversacion tecnica y comercial con
Kaleido. \texttt{FlowTwin} es un nombre provisional. La reunion no debe
vender un modelo ya validado: debe conseguir acceso controlado a
evidencia suficiente para decidir si existe un producto rentable.

\subsection{La idea en una frase}\label{la-idea-en-una-frase}

Convertir los eventos que Kaleido ya registra en Trace Port en una capa
predictiva que estime tiempo restante y riesgo de desviacion, explique
donde se forma el cuello de botella y compare escenarios operativos
aprobados, sin tomar el control de la operacion.

\subsection{El mensaje que debe
quedar}\label{el-mensaje-que-debe-quedar}

Kaleido ya tiene la parte dificil que muchos proyectos de IA no tienen:
un producto conectado al trabajo real, eventos con significado operativo
y acceso a usuarios que pueden actuar. No proponemos sustituir Trace
Port ni construir un simulador visual. Proponemos convertir esa
trazabilidad en anticipacion.

La falta de un historico rico no se oculta ni se discute: es la primera
restriccion de diseno. La fase inicial mide la calidad y densidad de la
evidencia y entrega process mining, baselines y un simulador discreto
aun cuando no haya datos suficientes para entrenar un modelo profundo.

\subsection{Apertura de 60 segundos}\label{apertura-de-60-segundos}

\begin{quote}
Hemos revisado vuestra observacion sobre el historico y hemos cambiado
el planteamiento para que el proyecto no dependa de promesas de deep
learning. Trace Port ya registra proyectos, turnos, eventos, equipos,
incidencias y tiempos. Nuestra propuesta es empezar por una sola
operacion repetible y usar esos eventos para reconstruir el proceso
real, medir desviaciones y crear un baseline de tiempo restante y riesgo
a 2, 4 y 8 horas. Solo si un modelo secuencial o JEPA mejora ese
baseline fuera de muestra se incorpora. El resultado seria una capa
read-only sobre Trace Port: alerta anticipada, intervalo de confianza,
cuello de botella y escenarios operativos aprobados. En cuatro a seis
semanas no prometemos un gemelo autonomo; prometemos una decision
medible sobre producto, datos y retorno.
\end{quote}

\subsection{Version de cinco minutos}\label{version-de-cinco-minutos}

\begin{enumerate}
\def\labelenumi{\arabic{enumi}.}
\tightlist
\item
  \textbf{Problema.} Un dashboard explica que ha pasado; el valor
  adicional aparece cuando permite saber con antelacion si una operacion
  no va a cumplir el plan.
\item
  \textbf{Activo de Kaleido.} Trace Port estructura la actividad que
  necesitamos: proyecto, packing list, turno, evento, equipo, evidencia
  e incidencia.
\item
  \textbf{Primer caso.} En una carga, descarga o manipulacion repetible,
  estimar tiempo restante y probabilidad de exceder el plan en
  horizontes de 2/4/8 h.
\item
  \textbf{Salida.} Riesgo calibrado, intervalo, factores asociados,
  cuello de botella, abstencion si falta evidencia y comparacion de
  acciones permitidas.
\item
  \textbf{Metodo.} Auditoria temporal, process mining, modelos simples,
  simulacion de eventos discretos y, como candidato, representacion
  secuencial tipo JEPA.
\item
  \textbf{Prueba honesta.} Separacion por proyecto y tiempo; nada del
  futuro entra en el pasado; se compara con reglas y modelos tabulares;
  se miden falsas alarmas, calibracion y minutos de anticipacion.
\item
  \textbf{Negocio.} Modulo premium de Trace Port, servicio de
  implantacion y mejora interna de planificacion. El precio y el ROI se
  calculan con el coste real de una desviacion y de intervenir, no con
  porcentajes genericos.
\item
  \textbf{Siguiente paso.} Una sesion de esquema y tres a cinco
  operaciones anonimizadas para producir un informe de viabilidad y un
  contrato de datos.
\end{enumerate}

\subsection{Objetivo concreto de la
reunion}\label{objetivo-concreto-de-la-reunion}

Salir con cuatro decisiones:

\begin{itemize}
\tightlist
\item
  propietario de negocio y propietario tecnico;
\item
  una operacion repetible con una decision que todavia pueda cambiarse;
\item
  acceso a esquema y muestra anonimizada, preferiblemente tres a cinco
  casos;
\item
  fecha de una sesion de 90 minutos para mapear proceso, costes y
  acciones.
\end{itemize}

No hace falta conseguir en la primera reunion acceso a produccion,
presupuesto cerrado ni compromiso con JEPA.

\subsection{Descubrimiento: las preguntas que
importan}\label{descubrimiento-las-preguntas-que-importan}

\subsubsection{Operacion y decision}\label{operacion-y-decision}

\begin{itemize}
\tightlist
\item
  ¿Que operacion se repite con suficiente frecuencia y tiene principio y
  fin?
\item
  ¿Que plan o SLA se incumple y cuando deja de ser util advertirlo?
\item
  ¿Quien puede actuar ante una alerta y que acciones tiene permitidas?
\item
  ¿Cual es el coste aproximado de retraso, recurso ocioso, hora extra,
  demurrage, penalizacion o replanificacion?
\item
  ¿Que accion es segura pero cara? Ese dato determina el umbral de
  alerta.
\end{itemize}

\subsubsection{Evidencia disponible}\label{evidencia-disponible}

\begin{itemize}
\tightlist
\item
  ¿Cada evento conserva \texttt{event\_time}, \texttt{ingest\_time},
  proyecto, turno y actor?
\item
  ¿Se guarda el plan que estaba vigente en ese momento o solo su version
  final?
\item
  ¿Hay reintentos, correcciones, eventos tardios, duplicados o trabajo
  offline?
\item
  ¿Se conocen pausas, cambios de equipo y causa de incidencia?
\item
  ¿Que parte existe en Trace Port y cual vive en ERP, TOS, Excel o
  correo?
\item
  ¿Cuantos proyectos completos hay por tipo de operacion y por terminal?
\end{itemize}

\subsubsection{Producto e integracion}\label{producto-e-integracion}

\begin{itemize}
\tightlist
\item
  ¿El mejor punto de entrada es API, export programado o replica
  read-only?
\item
  ¿Que interfaz consume el usuario: panel Trace Port, email, push o API?
\item
  ¿Debe desplegarse en infraestructura de Kaleido, nube privada o
  entorno del cliente final?
\item
  ¿Quien comercializaria el modulo y como se factura hoy Trace Port?
\end{itemize}

\subsection{Respuesta a la objecion principal: ``no hay historico
rico''}\label{respuesta-a-la-objecion-principal-no-hay-historico-rico}

Respuesta breve:

\begin{quote}
Precisamente por eso proponemos un piloto con salida util antes del
modelo. La primera entrega mide el proceso, la calidad temporal y que
señal existe. Con pocos datos usamos reglas, pooling por familias de
operacion, simulacion y estimadores con intervalos amplios. Si la
evidencia no soporta prediccion, el resultado es un no-go fundamentado y
un plan de instrumentacion, no una demo que inventa precision.
\end{quote}

Consecuencias tecnicas:

\begin{itemize}
\tightlist
\item
  el historico se mide en casos completos y densidad de eventos, no solo
  meses;
\item
  se agrupan operaciones solo si comparten mecanismo y contrato
  semantico;
\item
  el modelo se abstiene cuando el caso esta fuera de distribucion;
\item
  un log publico o sintetico valida tuberias, nunca el valor para
  Kaleido;
\item
  los datos nuevos se capturan con un contrato que permita aprendizaje
  futuro.
\end{itemize}

\subsection{Que significa aqui ``world
model''}\label{que-significa-aqui-world-model}

No es un generador de video ni una maqueta 3D. Es un modelo del estado
operativo que intenta responder:

\begin{enumerate}
\def\labelenumi{\arabic{enumi}.}
\tightlist
\item
  donde esta ahora la operacion;
\item
  que puede ocurrir a continuacion;
\item
  con que incertidumbre y en que plazo;
\item
  como cambia la distribucion si se modifica una accion permitida.
\end{enumerate}

Formalmente, para estado historico \texttt{h\_t}, contexto \texttt{c\_t}
y accion candidata \texttt{a\_t}, se estima una distribucion futura:

\texttt{p(y\_\{t+h\},\ z\_\{t+h\}\ \textbar{}\ h\_t,\ c\_t,\ a\_t)}

\texttt{z} es un estado latente; \texttt{y} contiene tiempos, eventos e
incidencias observables. Sin variacion real o una simulacion validada de
\texttt{a\_t}, el sistema solo puede mostrar asociacion o escenario, no
efecto causal.

\subsection{Por que JEPA, y por que no es un
requisito}\label{por-que-jepa-y-por-que-no-es-un-requisito}

JEPA aprende a predecir representaciones de partes futuras u ocultas,
evitando reconstruir cada detalle del dato. Puede ser util cuando hay
muchos eventos sin etiquetas limpias, distintas cadencias y estados
parcialmente observados. Para Kaleido el candidato seria un Event-JEPA
temporal, no V-JEPA de video.

No obstante, el modelo ganador puede ser Kaplan-Meier, gradient boosting
o una red secuencial pequena. JEPA entra solo si aporta una mejora
estable en holdout, calibracion y anticipacion a coste operativo
comparable. La arquitectura de producto no depende de esa eleccion.

\subsection{Arquitectura que conviene explicar en la
pizarra}\label{arquitectura-que-conviene-explicar-en-la-pizarra}

\begin{Shaded}
\begin{Highlighting}[]
\NormalTok{Trace Port / ERP / TOS / meteo}
\NormalTok{              |}
\NormalTok{              v}
\NormalTok{     eventos versionados + calidad temporal}
\NormalTok{              |}
\NormalTok{       estado operativo a tiempo t}
\NormalTok{              |}
\NormalTok{   +{-}{-}{-}{-}{-}{-}{-}{-}{-}{-}+{-}{-}{-}{-}{-}{-}{-}{-}{-}{-}+}
\NormalTok{   |          |          |}
\NormalTok{baseline   Event{-}JEPA   simulacion discreta}
\NormalTok{   |          |          |}
\NormalTok{   +{-}{-}{-}{-}{-}{-}\textgreater{} riesgo, tiempo, intervalo}
\NormalTok{                    |}
\NormalTok{          politica read{-}only + abstencion}
\NormalTok{                    |}
\NormalTok{      Trace Port / API / informe de turno}
\end{Highlighting}
\end{Shaded}

Puntos tecnicos que generan confianza:

\begin{itemize}
\tightlist
\item
  \texttt{event\_time} y \texttt{ingest\_time} se conservan por
  separado;
\item
  los planes se versionan para no usar el plan final al evaluar el
  pasado;
\item
  el split es por proyecto/tiempo/terminal, nunca por filas aleatorias;
\item
  la prediccion se hace como si solo se supiera lo disponible en ese
  instante;
\item
  cada salida conserva modelo, version de datos y explicacion;
\item
  no hay escritura automatica sobre sistemas operativos o de seguridad.
\end{itemize}

\subsection{Prueba y criterios de
salida}\label{prueba-y-criterios-de-salida}

El baseline no es ``acertar la media''. Incluye reglas operativas,
mediana por familia, survival analysis y boosting tabular. El candidato
debe probar valor incremental en un holdout congelado.

Metricas necesarias:

\begin{itemize}
\tightlist
\item
  MAE y pinball loss del tiempo restante;
\item
  AUCPR para eventos raros, no solo accuracy o AUROC;
\item
  Brier score y error de calibracion;
\item
  lead time antes de la desviacion;
\item
  falsas alertas por operacion o turno;
\item
  cobertura y anchura de intervalos;
\item
  rendimiento por tipo de carga, terminal y turno;
\item
  utilidad economica a distintos costes de intervenir y no intervenir.
\end{itemize}

Gates:

\begin{itemize}
\tightlist
\item
  \textbf{go:} mejora repetible y alerta con tiempo accionable;
\item
  \textbf{instrument:} el proceso tiene valor, pero faltan campos o
  consistencia;
\item
  \textbf{no-go:} no hay una decision accionable, no se puede evitar
  fuga temporal o el modelo no supera reglas simples.
\end{itemize}

\subsection{Encaje con los productos de
Kaleido}\label{encaje-con-los-productos-de-kaleido}

\subsubsection{Trace Port: primer
modulo}\label{trace-port-primer-modulo}

``Predictive Operations'' sobre la pantalla de proyecto/turno: tiempo
restante, riesgo, confianza, cuello de botella y enlace a evidencia. Es
el mejor lugar de inicio porque ya dispone de eventos operativos y
usuario en contexto.

\subsubsection{Shipping Board: segunda
extension}\label{shipping-board-segunda-extension}

Riesgo de retraso o excepcion en expedicion/recepcion, con
sincronizacion de documentos y eventos externos. Se reutiliza el motor,
pero cambia el contrato de entidades y horizontes.

\subsubsection{Freight Intelligence: tercera
extension}\label{freight-intelligence-tercera-extension}

ETA probabilistica y riesgo de excepcion de contenedor. Requiere separar
eventos del carrier, estimaciones externas y conocimiento disponible en
cada instante.

\subsection{Hipotesis de negocio}\label{hipotesis-de-negocio}

No presentar una cifra de precio cerrada sin conocer empaquetado ni
costes. Si Kaleido confirma valor, hay cuatro vias compatibles:

\begin{itemize}
\tightlist
\item
  add-on SaaS por terminal, proyecto o volumen de operaciones;
\item
  implantacion y calibracion inicial por cliente;
\item
  white-label dentro de Trace Port;
\item
  ahorro interno en planificacion, supervision y revision de operativas.
\end{itemize}

Formula para construir el caso:

\texttt{valor\ esperado\ =\ desviaciones\ evitadas\ +\ horas\ ahorradas\ +\ mejor\ uso\ de\ recursos\ -\ coste\ de\ intervenciones\ -\ coste\ del\ modulo}

El piloto debe registrar decisiones y costes para calcular esta
expresion con datos de Kaleido.

\subsection{Objeciones y respuestas}\label{objeciones-y-respuestas}

\subsubsection{``Ya tenemos dashboards y
alertas''}\label{ya-tenemos-dashboards-y-alertas}

FlowTwin no duplica visibilidad. La pregunta es si la alerta llega
antes, esta calibrada y permite una accion que reduzca coste. Si una
regla existente lo hace igual de bien, esa regla es el producto
correcto.

\subsubsection{``¿Esto es un gemelo
digital?''}\label{esto-es-un-gemelo-digital}

Es un gemelo operativo ligero: estado, transicion, incertidumbre y
escenarios. No pretende reproducir fisica, geometria o video del puerto
en la primera fase.

\subsubsection{``¿La recomendacion es
causal?''}\label{la-recomendacion-es-causal}

No por defecto. Sin datos de acciones, aleatorizacion o una simulacion
validada, se etiquetara como asociacion o escenario. Las decisiones de
alto impacto siguen siendo humanas.

\subsubsection{``¿Por que no usar directamente un
LLM?''}\label{por-que-no-usar-directamente-un-llm}

Un LLM puede resumir informes o explicar resultados, pero no sustituye
un modelo temporal calibrado ni soluciona la fuga de informacion. Puede
ser interfaz, no fuente de verdad numerica.

\subsubsection{``¿Necesitamos enviar datos
sensibles?''}\label{necesitamos-enviar-datos-sensibles}

No necesariamente. Se puede empezar con esquema, perfiles agregados y
muestra pseudonimizada; despues desplegar cerca del dato. Se eliminan
nombres y texto libre no necesario, y se acuerdan retencion, roles y
finalidad.

\subsubsection{``¿Cuanto tarda?''}\label{cuanto-tarda}

Cuatro a seis semanas para llegar a una decision de producto si el
acceso y el responsable operativo estan disponibles. No equivale a
integracion de produccion ni garantiza que un modelo avanzado supere el
baseline.

\subsection{Agenda recomendada de 60
minutos}\label{agenda-recomendada-de-60-minutos}

{\def\LTcaptype{none} % do not increment counter
\begin{longtable}[]{@{}rll@{}}
\toprule\noalign{}
Minutos & Tema & Resultado \\
\midrule\noalign{}
\endhead
\bottomrule\noalign{}
\endlastfoot
0-5 & Objetivo y restriccion de datos & Alineamiento \\
5-15 & Demo o recorrido de una operacion & Flujo real \\
15-30 & Eventos, planes, incidencias y sistemas & Mapa de evidencia \\
30-40 & Decision, acciones y coste & Target accionable \\
40-50 & Propuesta y gates & Alcance del scan/piloto \\
50-60 & Propietarios, muestra y fechas & Siguiente paso \\
\end{longtable}
}

\subsection{Que mostrar y que no
mostrar}\label{que-mostrar-y-que-no-mostrar}

Mostrar:

\begin{itemize}
\tightlist
\item
  el diagrama de la arquitectura y la escalera de cold start;
\item
  un ejemplo ficticio claramente marcado como tal;
\item
  criterios de go/no-go y peticion minima de datos;
\item
  que el sistema complementa los productos actuales de Kaleido.
\end{itemize}

No mostrar como evidencia:

\begin{itemize}
\tightlist
\item
  metricas antiguas de \texttt{industrial\_jepa\_mvp},
  \texttt{e-jepa-ttc} o \texttt{softnav-jepa};
\item
  una precision sin prevalencia, holdout y baseline;
\item
  resultados de datasets publicos como prueba de ROI para Kaleido;
\item
  ``causal'', ``autonomo'', ``SOTA'' o ``gemelo completo'' sin prueba.
\end{itemize}

\subsection{Lectura de los repositorios
existentes}\label{lectura-de-los-repositorios-existentes}

\begin{itemize}
\tightlist
\item
  \texttt{predictiveops-worldmodel} aporta contratos, manifests, splits,
  tests y gates.
\item
  \texttt{industrial\_jepa\_mvp} aporta patrones de agregacion y
  reporting, no evidencia transferible a logistica.
\item
  \texttt{e-jepa-ttc} aporta ventanas temporales, horizontes y
  streaming.
\item
  \texttt{softnav-jepa} aporta la separacion entre estimador y politica
  segura.
\item
  \texttt{SemanticSegmentation3Dclouds} queda como opcion futura si
  aparece un caso LiDAR/3D, no para el MVP de eventos.
\end{itemize}

Ninguno es una aplicacion lista para Trace Port. Reutilizar
indiscriminadamente sus modelos introduciria supuestos de sensores,
labels y splits incompatibles.

\subsection{Cierre de la reunion}\label{cierre-de-la-reunion}

\begin{quote}
Para no pediros un proyecto de IA a ciegas, proponemos empezar por una
muestra pequena y una operacion concreta. En una semana os devolvemos el
mapa de datos, la definicion exacta del target, los riesgos y un diseno
de evaluacion. Con eso decidimos juntos si merece la pena ejecutar el
piloto de cuatro a seis semanas o si lo correcto es primero mejorar
instrumentacion.
\end{quote}

\subsection{Lista personal antes de
entrar}\label{lista-personal-antes-de-entrar}

\begin{itemize}
\tightlist
\item
  Llevar la presentacion PDF y la version HTML offline.
\item
  Tener abiertos \texttt{DATA\_REQUEST.md} y \texttt{ARCHITECTURE.md}.
\item
  Saber explicar diferencia entre evento, caso, objeto y snapshot.
\item
  Preguntar, no asumir, que Trace Port conserva historial de versiones
  del plan.
\item
  No discutir algoritmo antes de fijar decision y coste.
\item
  Anotar vocabulario exacto de Kaleido para estados, incidencias y
  recursos.
\item
  Terminar con una peticion pequena y fechada.
\end{itemize}

\subsection{Glosario rapido}\label{glosario-rapido}

\begin{itemize}
\tightlist
\item
  \textbf{Caso:} una ejecucion completa que puede evaluarse, por ejemplo
  una operacion.
\item
  \textbf{Objeto:} entidad que participa en varios eventos, como
  proyecto, turno, equipo, buque o unidad de carga.
\item
  \textbf{OCEL:} log de eventos centrado en objetos; evita forzar todo a
  un unico caso.
\item
  \textbf{Process mining:} reconstruccion y comparacion del proceso real
  desde eventos.
\item
  \textbf{Conformance:} medida de divergencia entre proceso observado y
  esperado.
\item
  \textbf{World model:} modelo de estado y futuros posibles
  condicionados al contexto.
\item
  \textbf{JEPA:} prediccion en espacio de representaciones, no de cada
  detalle crudo.
\item
  \textbf{Calibracion:} que un riesgo del 20 \% ocurra aproximadamente
  dos de cada diez veces en poblaciones comparables.
\item
  \textbf{Lead time:} anticipacion entre alerta y evento objetivo.
\item
  \textbf{Abstencion:} no emitir una recomendacion cuando la evidencia
  no es fiable.
\item
  \textbf{Holdout:} conjunto de evaluacion nunca usado para seleccionar
  el modelo.
\end{itemize}

\end{document}
